\setlength{\parindent}{20pt}

\chapter{Sobre el dolç fluid}

\epigraf{Aquest noi és un pervertit.}{Emili Bagan}

\lettrine[loversize=0.2, lines=2]{A}{companyat} de Transímaco i Niní, nets de Ioulía, vaig baixar ahir als jardins del Tauleo, mogut pel propòsit d'orar a la deessa Maola, i amb ganes de veure la festa que s'hi trobà per allà, per primer cop donat que la celebraven per primera vegada aquella primavera.

Un cop havíem acabat les nostres oracions i contemplat l’espectacle, ens disposàvem a tornar a la ciutat, quan Rógos, fill de Phrangískos, ens veié de lluny i digué al seu servent, Sérgēs, que corregués a demanar-nos que l’esperéssim.

\begin{platdialogue}
\item[\speaker{Sérgēs:}] Rógos us demana que espereu.
\end{platdialogue}
Girant-me jo aleshores l'hi preguntè on era ell.
\begin{platdialogue}
\item[\speaker{Sérgēs:}] Helo allà enrere --- contestà --- ara s'apropà, espereue'l
\item[\speaker{Dauídas:}] Bè, sigui doncs així, l'esperarem.
\end{platdialogue}
En efecte, poc després arribaren Rógos, amb Dareíos, el germà d'Aléxandros, Karlótta i Einará i alguns més, probablement provenents de la festa.
\begin{platdialogue}
\item[\speaker{Rógos:}] Segons el meu judici, sembla que, Dauídas marxeu-vos ja de tornada a la cuitat. 
\item[\speaker{Dauídas:}] I no t'equivoques.
\item[\speaker{Rógos:}] Veus, doncs, quants som nosaltres?
\item[\speaker{Dauídas:}] És resposta vàlida qualsevol altra que trivial?
\item[\speaker{Rógos:}] Així doncs, o teniu la força per vèncer-nos a tots, o us romandreu aquí.
\item[\speaker{Dauídas:}] I no hi ha doncs, una altra sortida?
\item[\speaker{Dareiós:}] De cap manera.
\item[\speaker{Karlótta:}] No sou potser coneixedors de què celebrarem una carrera de cavalls amb torxes al vespre en honor a la deessa?
\item[\speaker{Dareiós:}] A cavall?! A la taula aquella, com pot ser doncs?!
\item[\speaker{Rógos:}] Com ho sents.
\item[\speaker{Niní:}] Veig doncs que ens haurem de quedar.
\item[\speaker{Dauídas:}] Doncs això sembla.
\end{platdialogue}
Ens vam dirigir doncs, a la tataula, casa de Rógos i ens trobarem allà una bona festa muntada. Estava allà Aimílios, aví de Rógos, que em semblà molt avançat en anys, doncs feia temps que no el veia, ell tenia posada una corona, doncs havien realitzat un sacrifici al riu poc abans. Nosaltres ens vàrem asseure a la tataula, al seu voltant. Al veurem, Aimílios em saluda i digué:
\begin{platdialogue}
\item[\speaker{Aimílios:}] Oh Dauídas!, com són de rares les vegades que baixes a veure'ns al Tauleo. No hauria de ser així, doncs si tinguéssim la valentia de pujar a la C3 no faria falta que vinguessis aquí sinó que aniríem nosaltres al que ja és casa teva, però com no és així sou vosaltres els que heu de baixar aquí amb més freqüència.
\item[\speaker{Rógos:}] Dauídas --- va dir tallant al seu avi --- voldria conversar d'un tema amb vos si no tens inconvenient.
\item[\speaker{Dauídas:}] Com em podria negar a tal ofrenda. De què es tracta?
\item[\speaker{Rógos:}] Doncs mira, Oh Dauídas. L'altra vesprada, sota cartes i humor de civada varem generar el següent dubte: Amb quin propòsit s’acostà aquell primer home a una vaca per extreure’n llet?
\item[\speaker{Niní:}]Ho dius de debò? Això és digne d’un banquet, no d’un diàleg filosòfic.
\item[\speaker{Dauídas:}] Potser sí, Niní, però no tot allò digne de riure és indigne de ser examinat. Car no és pas el riure el que desautoritza una qüestió, sinó la manca de raó que l’acompanya. És doncs una gran qüestió. Oh Rógos, anem a desglossar-ho per veure quines conclusions n'extrèiem.
\item[\speaker{Rógos:}] Com no, Oh Dauídas!
\item[\speaker{Dauídas:}] Així doncs, primer plantegem-nos el primer pas. Aquell primer home, per què ho va fer?
\item[\speaker{Karlótta:}] Potser tenia gana i prou. La gana és mestra d’astúcia. I, potser, en veure aquell animal robust, amb ventre ple i flancs tibants, pressuposà que allà hi havia nodriment amagat, com un tresor natural esperant ser descobert. Així, empès per la necessitat, però guiat per una suposició, s’hi acostà amb ànsia mesclada d’esperança.
\item[\speaker{Dauídas:}] És cert, oh Karlótta, però com sabia l'home que allà hi havia menjar, no em sembla, en primera vista, que, en aquell racó de la natura, s'hi amagui tal tresor. L'home sens dubte, sense cap certesa, segurament empès per la necessitat, però guiat per una suposició, s’hi acostà amb ànsia mesclada d’esperança. I potser —no ho sabem— no cercava només omplir l’estómac, sinó que hi havia en ell un desig més fosc o més profund. Car, qui pot dir què mou realment l’home quan posa la mà allà on la naturalesa ha posat la font de la vida? ¿Era la recerca d’aliment, o també una mena de curiositat sacra, o fins i tot un impuls que barrejava el desig amb el misteri?
\item[\speaker{Karlótta:}] Dauídas, vols insinuar que potser hi havia en aquell gest un rastre d’erotisme?
\item[\speaker{Rógos:}] No ho descartaria pas. El cos de la vaca, amb les seves corbes generoses, no és pas neutre per als sentits. ¿No podria ser que el primer gest de munyir portés també una confusió entre el plaer i la utilitat? L’home primitiu encara no separava allò sagrat, allò útil i allò desitjable.
\item[\speaker{Dauídas:}] Podria ben ser. --- Vaig comentar--- però crec que mai podrem conèixer les vertaderes intencions, el que sí que podem saber de bon grat és que molt probablement aquell home no tenia clar quines eren les seves intencions, i si fos així, això revelaria quelcom encara més profund: que el coneixement neix sovint d’una barreja d’impulsos, no pas de puresa. Que el camí cap a la veritat no sempre comença amb una pregunta clara, sinó amb una agitació obscura del cos i de l’esperit.
\item[\speaker{Rógos:}] Crec que tens raó, Oh Dauídas!
\item[\speaker{Dauídas:}] Així doncs, si seguim la idea que l’home estava guiat per una intuïció més enllà de la mera necessitat, podem pensar que sovint el cos actua abans que l’ànima comprengui allò que realment desitja. I pot ser que, tocant el ventre de la vaca, aquell home cercava una forma de retrobar la font primigènia de l’aliment matern, una mena de retorn a la mare perduda de la infància de la humanitat.
\item[\speaker{Rógos:}] Això ja no és gana, Dauídas, sinó nostàlgia.
\item[\speaker{Dauídas:}] I potser la nostàlgia és una de les formes més profundes de coneixement, Rógos. Però ¿la necessitat sola pot conduir a l’enginy? ¿O cal també un cert enteniment, o almenys una esperança que allò desconegut pugui ser profitós?
\item[\speaker{Karlótta:}]Potser ho va veure fer a un animal, i va pensar que, si era bo per al vedell, ho seria també per a ell.
\item[\speaker{Dauídas:}] Aleshores dius que l’home aprèn per imitació. Però digues, ¿és la imitació una forma de coneixement veritable, o només una ombra de la veritat? ¿Com distingim entre la còpia i l’essència?
\item[\speaker{Rógos:}] Sempre acabes portant-nos cap a les idees...
\item[\speaker{Dauídas:}]I on hauríem d’anar, si no és cap al que és en si? ¿O preferiu que restem entre les coses que només semblen, com si fóssim homes en una cova, prenent ombres per realitats?
\end{platdialogue}
Interromp Dareiós, amb actitud serena i mirada oberta.
\begin{platdialogue}
\item[\speaker{Dareiós:}] Us sento parlar de munyir vaques i em sembla que no parleu de ramaderia, sinó d’alguna altra cosa més elevada. Quan jo era jove, a Dinamarca, els homes munyien per tradició. Cap de nosaltres no sabia d’on havia sortit aquell costum. Però ningú no ho qüestionava.
\item[\speaker{Dauídas:}] Així doncs, el costum pren el lloc del coneixement. Però no és costum també el que perpetua l’error? Digueu-me, ¿com podem saber si el primer munyidor era un savi, un foll, o un inspirat pels déus?
\item[\speaker{Rógos:}] Potser era totes tres coses a la vegada.
\item[\speaker{Dauídas:}] I això ens du a una altra qüestió, potser encara més profunda: ¿què és, en realitat, la saviesa? ¿No serà, potser, la capacitat d’orientar-se en la foscor, d’estirar el fil d’allò desconegut fins a trobar-ne la utilitat?
\item[\speaker{Rógos:}] I si és així, el primer munyidor seria un filòsof.
\item[\speaker{Dauídas:}] O si més no, un aspirant a la filosofia, que és l’amor a allò que encara no es posseeix, però que es desitja amb ànima sincera.
\item[\speaker{Dareiós:}] Doncs jo crec que estava tocat de l'ala.
\item[\speaker{Dauídas:}] No ho sabrem mai, no avançarem mai de l'especulació, però passant a la següent qüestió, l'home, en probar el líquid i haver descobert el tresor es va quedar allà?
\item[\speaker{Rógos:}] Que vols dir Oh Dauídas?!
\item[\speaker{Dauídas:}] Un cop obtinguda tal recompensa d’un animal determinat, és raonable suposar que l’home s’aventurà a provar amb d’altres. I així, potser, nasqué el costum d’extreure llet de la cabra o de l’ovella, bèsties més petites i diferents de la vaca. Però digueu-me, amics: ¿no és també versemblant que, després d’haver trobat fruit en la femella, l’home, desconeixent encara la naturalesa d’aquell do, intentés el mateix amb el mascle?
\item[\speaker{Karlótta:}] Molt bona no hagues set la recompensa.
\item[\speaker{Rógos:}] Debatible, Oh Karlótta!
\item[\speaker{Dauídas:}] I cal considerar, amics meus, que un cop un cert saber ha estat establert, l’home no s’atura pas en la seva aplicació original. Ans al contrari, impulsat per l’èxit, s’afanya a estendre aquell coneixement a d’altres àmbits, com qui prova si una clau obre no pas una sola porta, sinó totes les estances de la casa. Així, després d’haver trobat en la femella una font de nodriment, no és difícil imaginar que ell, o potser els seus descendents, temptessin fortuna també en altres animals, moguts no per coneixement cert, sinó per la voluntat de saber fins on arriba la nova llei descoberta. ¿I no és pròpia de l’ànima humana aquesta impaciència per universalitzar allò que ha començat com un sol encert? Fins on va arribar aquest home?
\item[\speaker{Rógos:}] Avui coneixem més de deu bèsties de les quals l’home ha sabut extreure llet de manera profitosa. I si tants han estat els èxits que ens han arribat, ¿quants errors degué cometre aquell primer esperit inquisitiu abans no trobà el camí segur?
\item[\speaker{Dauídas:}] Així és. Car si la naturalesa no mostra els seus fruits de bon principi, l’home, desconeixent encara les causes, actua per assaig i esperança. I potser va intentar munyir el cavall, el porc, el gos, fins i tot bèsties salvatges o mascles infructuosos, sense saber si allò que havia descobert era una regla universal o una excepció beneïda. En molts casos no va poder extreure cap brebatge, mentres que en altres extraient-ne un de sabor no beneficiós, com podria haver set el cas dels mascles.
\item[\speaker{Rógos:}] Segons el nostre raonament, doncs, és gairebé inevitable pensar que l’home, empès per l’èxit obtingut, acabaria provant també el licor del toro, el del moltó, i així successivament, en una llarga desfilada d’assaigs que abraçava el món animal sense distinció de sexe ni natura.
\item[\speaker{Dauídas:}] De ben segur.
\item[\speaker{Karlótta:}] I no és aquesta, Dauídas, la imatge mateixa de la condició humana? La d’aquell que camina a les palpentes, i ensopega més sovint del que encerta, però que es nega a restar quiet en la ignorància?
\item[\speaker{Dauídas:}]Ho has dit amb saviesa, Karlótta. Perquè la veritable ciència no neix només del descobriment, sinó també del reconeixement de l’error. I potser, més que els èxits, són els fracassos els que indiquen la direcció del coneixement. Així, cada animal que no respongué amb llet fou també mestre, per via negativa, de la veritat.
\end{platdialogue}